\documentclass[11pt]{article}
\usepackage[margin=1in]{geometry}
\usepackage{amsmath,amssymb,booktabs,array}
\usepackage[colorlinks=true,linkcolor=black,urlcolor=blue!55!black]{hyperref}
\usepackage{parskip}
\setlength{\parskip}{6pt}
\newcommand{\code}[1]{\texttt{\small #1}}

\title{\vspace{-1.2cm}\textbf{A simulator for adaptive tutoring}\\[3pt]
\large Student models, knowledge tracing, question selection, and what we
can and cannot measure with them}
\date{\small 27 September 2026}
\author{}

\begin{document}
\maketitle
\vspace{-1.1cm}
\noindent\rule{\textwidth}{0.4pt}

\section*{1. What the simulator is for}

A tutoring system has to choose the next question without being able to see
what the student knows. To study that choice we need a student whose
knowledge we \emph{can} see, so that any claim about the tutor can be checked
against the truth rather than against another estimate. Everything here is
built around that one idea: the simulator owns a hidden state, the tutor
never sees it, and every measurement compares the two.

The simulator has three replaceable parts, and keeping them separate is what
makes the measurements interpretable.

\begin{center}
\small
\begin{tabular}{@{}llp{7.4cm}@{}}
\toprule
\textbf{Part} & \textbf{Sees the truth?} & \textbf{Question it answers}\\
\midrule
Student model & it \emph{is} the truth & what does the learner actually know, and how does that change?\\
Knowledge tracing & no & what does the tutor believe, given only right and wrong answers?\\
Question selection & no & given that belief, what should be asked next?\\
\bottomrule
\end{tabular}
\end{center}

Because they are separate, one can be varied with the others held fixed. That
is the only way to attribute an effect: if the student and the tracer are
fixed and only the scheduler changes, any difference is the scheduler's.

\section*{2. The curriculum}

A curriculum is a directed acyclic graph of concepts. Each concept carries a
tier (its depth), a list of prerequisites, and a bank of practice questions.
Curricula are data, stored as JSON, so a new one is a new file rather than a
code change.

The default graph is deliberately abstract and small: eight concepts named by
tier, \code{A1, A2} at the top with no prerequisites, then \code{B1, B2, B3},
then \code{C1, C2}, and \code{D1} as the single capstone. A name states its
own depth. The shape is a symmetric diamond, and two properties of it matter.
The zone of proximal development starts at exactly two concepts, small enough
to read at a glance; and \code{B2} is a prerequisite of both \code{C1} and
\code{C2}, so a single weak concept visibly blocks two downstream. A
40-concept language curriculum is kept alongside it so that earlier results
stay reproducible.

\section*{3. The student models}

\subsection*{3.1 The plain BKT student}

Each concept is a hidden coin: known, or not. Four parameters govern it.

\begin{center}
\small
\begin{tabular}{@{}cll@{}}
\toprule
 & \textbf{Meaning} & \textbf{Default}\\
\midrule
$p(L_0)$ & probability the concept is known before any question & 0.15\\
$p(T)$   & probability an unknown concept becomes known at an attempt & 0.12\\
$p(G)$   & probability of a correct answer while it is unknown & 0.25\\
$p(S)$   & probability of a wrong answer while it is known & 0.10\\
\bottomrule
\end{tabular}
\end{center}

At each attempt on concept $c$, with $K_c\in\{0,1\}$ the hidden state:
\begin{equation}
\Pr(\text{correct}) =
\begin{cases}
1-p(S), & K_c = 1,\\
p(G),   & K_c = 0,
\end{cases}
\qquad\text{then}\qquad
K_c \leftarrow 1 \ \text{ with probability } p(T) \text{ if } K_c = 0 .
\end{equation}

\textbf{The student answers from the state it was in, and only then
transitions.} A learning event therefore shows up on the \emph{next} question,
not the one that caused it, which is why a student can be seen learning a
concept on a question they got wrong. This ordering is not cosmetic: it is the
ordering the tutor's update assumes, and getting it backwards silently breaks
calibration.

There is no forgetting, no question difficulty, and by default no effect of
prerequisites on learning. That is the point. It is the simplest student for
which the tutor's arithmetic has a known right answer.

\subsection*{3.2 Two student models: S1 and S2}

\begin{description}
\item[S1] Prerequisites do not affect learning. A concept is learned at
$p(T)=0.2$ whenever it is asked, whatever else the student knows.
\item[S2] A concept is learned at $p(T)=0.2$ once \emph{all} of its
prerequisites are truly known, and at $p(T)=0.05$ otherwise.
\end{description}

S2 is the smallest change that makes the order of questions matter. Under S1
a question is worth the same whenever it is asked, so no scheduler can beat
any other by more than noise; under S2 a question asked too early is worth a
quarter as much.

\subsection*{3.3 The continuous student}

The original student is also still available: mastery is a real number in
$[0,1]$, questions carry difficulty, prerequisites gate the learning rate
continuously, and knowledge decays with elapsed time. It is the right model
for studying forgetting and spaced practice, and the wrong one for checking
the tutor's arithmetic, because it has no binary truth to compare a
probability against.

\section*{4. Knowledge tracing}

The tutor keeps one number per concept: $b_c$, its belief that the student
knows $c$. It sees only right and wrong answers. After an observation $X$ it
applies Bayes' rule and then the learning transition:
\begin{equation}
b^{+} =
\begin{cases}
\dfrac{b\,(1-p(S))}{b\,(1-p(S)) + (1-b)\,p(G)}, & X = \text{correct},\\[14pt]
\dfrac{b\,p(S)}{b\,p(S) + (1-b)\,(1-p(G))}, & X = \text{wrong},
\end{cases}
\qquad
b_c \leftarrow b^{+} + (1-b^{+})\,p(T).
\end{equation}

\textbf{Timing.} The belief written by this update is the belief about the
state the student is in \emph{going into the next question}, because the
transition has been applied. The student's hidden state after its own
transition refers to the same instant. Any comparison of belief against truth
must use that pairing; reading the state before the transition instead shifts
everything by one step and makes a correct tracer look badly calibrated. This
is the single easiest way to produce a convincing but wrong result, so it is
stated here and enforced in the tests.

Because the tutor's parameters can be set independently of the student's, the
tutor can be made deliberately wrong, which is what Section 7 exploits.

\section*{5. Question selection}

The zone of proximal development is the set of concepts the tutor believes are
worth teaching now:
\begin{equation}
\mathrm{ZPD}(t) = \{\,c : b_c < \theta \ \wedge\ b_p \ge \theta \ \ \forall p
\in \mathrm{prereq}(c)\,\},
\end{equation}
for a mastery threshold $\theta$. Note that it is defined entirely on
\emph{beliefs}: the tutor has no access to the student's real state.

Two rules are compared.

\begin{description}
\item[Q2, \code{uniform\_zpd}] uniform over the ZPD, so a concept is not asked
until its prerequisites are believed mastered. When nothing is
prerequisite-ready, which in this graph means everything is believed mastered,
it falls back to revisiting any concept.
\item[Q1, \code{uniform\_unmastered}] uniform over every concept not yet
believed mastered, ready or not.
\end{description}

Both are deliberately free of any bandit. The comparison is about
prerequisite gating alone.

\section*{6. What is measured, and against what}

Every metric reads the student's hidden state, never the tutor's belief.

\begin{description}
\item[Calibration.] Bin every belief the tutor has ever held; within a bin,
what fraction of the time was the student truly in the known state? A correct
tracer lies on the diagonal.
\item[Detection delay.] We know the exact step at which each concept was
learned. How many further questions \emph{on that concept} pass before the
tutor declares it mastered?
\item[False alarm rate.] Of all the declarations the tutor made, what share
came before the student had actually learned the concept?
\item[Concepts truly known.] How many concepts the student really knows at the
end. This is the only metric that can compare schedulers, because it owes
nothing to what any tutor believes.
\end{description}

Every comparison is paired: the same seeds are used in every arm, so learner
$i$ starts from the same hidden state in every arm and differences are read
within-learner rather than between groups. The pairing is on that shared
starting state, not on every coin thereafter: once two schedulers ask
different concepts the learners' random draws diverge, so the pairing buys
correlation rather than identity.

\section*{7. What the simulator has been used to establish}

\textbf{The tracer is correct.} On the plain BKT student with a correctly
specified tutor, pooled mean belief is 0.771 against a true known-fraction of
0.773 over 12{,}000 observations (100 learners, 120 questions each). At four
times the learners, 48{,}000 pairs, every bin holding at least 500
observations sits within 0.026 of the diagonal. Sweeping each of the four
parameters in turn over 19 settings leaves the pooled gap within 0.011
throughout. That last sweep keeps the tutor correctly specified in every cell,
so it shows the estimator is unbiased across parameter settings; it is not by
itself evidence that a \emph{mis}specified tutor would be caught.

\textbf{The two checks agree with each other.} If beliefs are calibrated, the
false-alarm rate at threshold $\theta$ must equal one minus the mean belief at
the moment of declaration. Measured and predicted agree at every threshold.
Two quantities computed by different routes landing on the same number is
stronger evidence than either alone.

\textbf{The scheduler only matters if the student model says it can.} Crossing
S1 and S2 with Q1 and Q2, 400 learners, 40 questions, scored on concepts the
student truly knows out of eight:

\begin{center}
\small
\begin{tabular}{@{}lcccc@{}}
\toprule
 & \textbf{Q1} & \textbf{Q2} & \textbf{Q2 $-$ Q1} & \textbf{95\% CI}\\
\midrule
S1 & 6.15 & 6.00 & $-0.15$ & $\pm 0.14$\\
S2 & 4.39 & 5.86 & $+1.47$ & $\pm 0.18$\\
\midrule
\multicolumn{3}{@{}l}{interaction (S2 $-$ S1)} & $+1.62$ & $\pm 0.17$\\
\bottomrule
\end{tabular}
\end{center}

\noindent The interaction is far from zero: the scheduler is worth about one
and a half concepts to S2 and nothing useful to S1. Because seed $i$ is the
same simulated person in all four cells, that interval is the paired one; the
two per-student differences are correlated ($r=+0.59$), so treating them as
independent would overstate the uncertainty. The S1 figure is small and, with
three contrasts in the table, marginal: it should be read as ``at most a
slight cost'', not as an established effect.

\textbf{What the restriction is actually worth.} Q1 and Q2 differ in two ways
at once, and separating them changes the story. Running a matched pair in
which \emph{both} arms gate on the student's true state, so that only the
prerequisite restriction differs:

\begin{center}
\small
\begin{tabular}{@{}lccc@{}}
\toprule
 & \textbf{unrestricted} & \textbf{restricted} & \textbf{effect}\\
\midrule
S1 & 7.480 & 7.480 & $0.000 \pm 0.000$\\
S2 & 5.513 & 7.480 & $+1.968 \pm 0.167$\\
\bottomrule
\end{tabular}
\end{center}

\noindent Under S1 the restriction is \emph{exactly} neutral: identical for
all 400 learners, not merely close. So the small negative figure in the table
above is not the cost of restricting to prerequisite-ready concepts. It is the
cost of deciding readiness from a \emph{belief that lags the truth}: Q2
withholds concepts the student is in fact already able to learn. The
whiteboard prediction that the restriction would be invisible under S1 was
right about the restriction; what it did not anticipate is the separate price
of gating on an estimate.

\textbf{Budget decides whether any of this is visible.} At 120 questions the
cells reach 7.99, 7.99, 7.68 and 7.90 of eight. The S1 effect is gone
($-0.003 \pm 0.015$), but the S2 effect is still there and still positive
($+0.212 \pm 0.094$). So a longer budget shrinks the contrast towards the
ceiling rather than abolishing it, and any comparison of question-selection
rules has to report the budget it was measured at.

\textbf{Misspecification degrades gracefully but not harmlessly.} Sweeping the
tutor's assumed $p(G)$ against a fixed student, 200 learners, with intervals
bootstrapped over learners: the pooled gap is $+0.002 \pm 0.006$ at the truth
and $-0.063 \pm 0.008$ when the tutor underestimates guessing by 0.20, while
false alarms move from 4.6\% to 17.1\%. With a population whose parameters are
spread uniformly by $\pm 0.15$ and a single fixed tutor, the gap is
$-0.017 \pm 0.014$ and false alarms rise from 4.6\% to 7.5\%. A single tutor
model stays usable across a varied population, at a measurable cost. Note the
pooled gap is a signed average, so it is a bias diagnostic: opposite-signed
bins can cancel, and per-bin behaviour is what Section 6's calibration check
looks at.

\section*{8. Limitations}

The plain BKT student has no forgetting, so nothing here speaks to retention or
spaced practice; the continuous student exists for that and is used in none of
the checks above. All results are on one curriculum of eight concepts, though
the interaction also reproduces on the 40-concept one. Every population in this
document derives from a single base seed, so the intervals describe sampling
within one seed block rather than across independent replications. Per-bin
calibration intervals are nominal: successive observations on one concept are
dependent, because the belief after one attempt is the belief before the next,
so the pooled comparison is the robust one. The threshold $\theta = 0.9$ is
shared between the tracer and the scheduler, which means it sets both when a
concept is declared and when its dependants become available, and it is not
swept. The tutor's $p(T)$ against S2 is held constant even though the student's
varies with its prerequisites, so the tutor is mildly misspecified there by
construction; whether it should instead gate $p(T)$ on its own belief about the
prerequisites is open. And every number is a statement about this model, not
about people: no result here has been checked against a real learner.

\vspace{6pt}
\noindent\rule{\textwidth}{0.4pt}
\vspace{2pt}

\noindent
\footnotesize
\textbf{Reproducing.} \code{python3 make\_sanity.py} regenerates every figure
and table quoted above into \code{sanity\_outputs/}; \code{python3
test\_sanity.py} and \code{python3 test\_simulator.py} run 29 and 55 checks
respectively, including the exactness checks that pin the update ordering
(with $p(G)=p(S)=0$ a correct answer must come from a known concept; adding
$p(T)=1$ forces the tutor's belief to equal the truth at every step). Seeds are
fixed throughout.

\end{document}
